\documentclass[a4paper]{article}
\usepackage[pdftex]{graphicx}
\usepackage[utf8]{inputenc}
\usepackage{enumerate}
\usepackage{siunitx}
\sisetup{locale=DE}
\usepackage{href-ul}
\hypersetup{
	colorlinks=true,
	linkcolor=blue,
	urlcolor=blue}
\usepackage{icomma}
\usepackage{amssymb}
\usepackage{geometry}
\geometry{a4paper, top=15mm, left=15mm, right=15mm, bottom=15mm,
	headsep=10mm, footskip=12mm}

\begin{document}
	{\bf Thermodynamics}
	\begin{enumerate}[1.]
		{\bf \item "Heat is the energy that is transferred between bodies as a result of temperature differences"} --- Briefly explain the three ways of transferring energy!
		
		
		%Exercise 2
		\begin{minipage}{0.5\textwidth}
			{\bf \item The diagram on the right shows a qualitative temperature curve of a process} in which heat was continuously supplied to a substance.
			\begin{enumerate}[a)]
				\item What substance could it be?
				\item Explain the terms melting, solidification, boiling, condensation, sublimation and resublimation.
				\item Explain what happens at curve points A to E
			\end{enumerate}
		\end{minipage}
		\hfill
		\begin{minipage}{0.45\textwidth}
			\includegraphics[width=5 cm]{wasser105.pdf}
		\end{minipage}
		
		{\bf \item 160 g of ice cream of $\SI{0}{\celsius}$ are in a container on a heating plate.} The ice cream should be completely converted into steam of $\SI{100}{\celsius}$. [$r_{H_2O} = \SI{2257}{\kilo\joule\over \kilogram};\quad s_{H_2O}=\SI{334}{\kilo\joule\over \kilogram};\quad c_{ H_2O}=\SI{4,18}{\kilo\joule\over \kilogram\kelvin}$]
		\begin{enumerate}[a)]
			\item What amount of heat must be added to the ice?
			\item How long does the heating plate have to give off the corresponding heat if it has a heating output of $P=\SI{2800}{\watt}$?
		\end{enumerate}
		
		
		\begin{minipage}{0.5\textwidth}
			{\bf \item The circuit diagram opposite shows a circuit with a voltage source, a load and a bimetal strip.}
			\begin{enumerate}
				\item Explain how a bimetallic strip works.
				\item Name this circuit and give an application of it.
			\end{enumerate}
		\end{minipage}
		\hfill
		\begin{minipage}{0.45\textwidth}
			\includegraphics[width=6 cm]{bimetall105.pdf}
		\end{minipage}
		
		
		
		{\bf \item Oxygen, which is present in the air in a proportion of approximately 21 \%, liquefies at approximately $ T=\SI{90}{\kelvin}$ above absolute zero.}\\
		By how many Kelvin do you have to increase the oxygen in the air at room temperature?
		$\SI{19}{\celsius}$ cool down so that it becomes liquid?
		
		
		{\bf \item $m_1=\SI{200}{\gram}$ Water vapor of temperature $\vartheta _{d} =\SI{100}{\celsius}$ are in $m_2=\SI{1800}{\ gram}$ water} of temperature $\vartheta _{w} = \SI{20}{\celsius} $ is introduced so that the water vapor condenses completely. The temperature of the liquid is then
		$\vartheta _{m} = \SI{82}{\celsius} $. The energy exchange with the environment and the calorimeter vessel should be neglected here.
		\begin{enumerate}[a)]
			\item What heat energy does the water of mass $m_2$ absorb?
			\item Calculate the specific condensation energy of water vapor!
		\end{enumerate}
		
		{\bf \item $m_w=3.2~kg$ Water of temperature $\vartheta_w = \SI{30}{\celsius}$ and $\SI{1.5}{\kilogram}$ ice} are in a calorimeter mixed. After the ice is 70\% melted, nothing changes. What temperature was the ice at the beginning?
		[Specific heat capacity of ice $c_{eis}= \SI{2,1}{\kilo\joule \over \kilogram\kelvin}$]
		
		
		%Task 1
		{\bf \item If you place a long metal spoon in a hot cup of coffee}, the part of the spoon protruding from the coffee also becomes hot.
		\begin{enumerate}[a)]
			\item Use the particle model to explain how this heat conduction works.
			\item Explain how the transfer of internal energy occurs through convection. Then explain this using an example of your own choice.
		\end{enumerate}
		
		%Exercise 2
		{\bf \item The aim is to calculate how much the drum brake} of a bicycle made of $\SI{1.0}{\kilogram}$ steel heats up when the brakes are constantly braked while driving downhill with an altitude difference of 100~m. The specific heat capacity of steel is $c=\SI{0.48}{\kilo\joule \over \kilogram\kelvin} $, the mass of the bicycle with rider is $\SI{85}{\kilogram}$.
			
			{\bf \item What is the first law of thermodynamics?}
		\end{enumerate}
			\fbox{
				\begin{minipage}{0.5\textwidth}
					For the solution, please \href{https://www.okuyakl.de/phys/p9thaL402/le402.pdf}{click here} or scan the QR code.\\
					Additional worksheets are available at
					
					\href{http://www.okuyakl.com}{www.okuyakl.com}
				\end{minipage}
				\hfill
				\begin{minipage}{0.4\textwidth}
					\includegraphics[width=1.5 cm]{../../viecher/zwe03}
					\includegraphics[width=3 cm]{qre402}
					\includegraphics[width=2 cm]{../../viecher/afanticon1}
					
			\end{minipage}}
			
	
	\end{document}